\chapter{Improper Integrals, Parameters, and the Gamma Function}
\label{ch:improper-parameters}

A specific improper integral needs bounds at its missing endpoints, not a
general theory of measurable functions.  A parameter derivative needs
corresponding bounds for its finite Taylor errors.  Gamma, beta, and the
error function give concrete tests of both requirements.

\section{Two tails and one bounded computation}
Suppose $f(t,z)$ has an integral computation on each truncated interval
$[\varepsilon,T]$, and uniform bounds on every finite end piece are supplied:
\[
 \left|\int_u^\varepsilon f(t,z)\,dt\right|\le e_0(\varepsilon),
 \quad
 \left|\int_T^V f(t,z)\,dt\right|\le e_\infty(T)
 \quad(0<u<\varepsilon<T<V).
\]
Both errors must shrink by proved schedules.  Evaluate the middle integral
to error $\eta$ and enlarge its rational box by
$e_0(\varepsilon)+e_\infty(T)$.  Comparisons on a larger common truncated
interval prove compatibility of any two such boxes.  Finite intersections
therefore construct $\int_0^\infty f(t,z)\,dt$.

These tail bounds can be absolute-integral bounds, but need not be: later
an oscillatory cancellation will give a bound without absolute convergence.
No interchange of operations is inferred merely from the notation
$\int_0^\infty$.

\section{Parameter series with an integrable majorant}
\label{thm:parameter-majorant}
Suppose on every truncated interval a supplied identity is
\[
 f(t,z_0+w)=\sum_{n\ge0}a_n(t)w^n,
 \qquad |a_n(t)|\le M(t)R^{-n},
\]
where $M$ is a nonnegative concrete integrand with proved endpoint tails
and an integral bounded by a rational $B$.  Each $a_n$ is integrable on
truncated intervals with the supplied continuity data.  The same bounds
construct $A_n=\int_0^\infty a_n(t)\,dt$.  For $|w|\le r<R$, finite
sums give
\[
 \left|\int_0^\infty f(t,z_0+w)\,dt-\sum_{n=0}^N A_nw^n\right|
 \le\frac{B(r/R)^{N+1}}{1-r/R},\qquad |A_n|\le BR^{-n}.
\]
To prove this, first work on $[\varepsilon,T]$, where everything is a
finite sum plus a uniform remainder.  Bound that remainder by the
geometric tail times $M(t)$.  The omitted end pieces have the same bound
with the corresponding tails of $M$.  Now make all three finite errors
small.  This proves termwise integration and constructs a local series
chart for the parameter function.  Its derivatives are the coefficient
computations from Chapter~\ref{ch:local-models}; no dominated-convergence
theorem is imported.

\section{Euler's gamma integral with uniform rational budgets}
For a rational complex $z$ with $\Re z>0$, define
\[
 \Gamma(z)=\int_0^\infty t^{z-1}E(-t)\,dt,
 \qquad t^{z-1}=E((z-1)L(t)).
\]
The same construction accepts computed inputs in a certified region
$a\le\Re z\le b$, where $a>0$ is rational.  On $0<t\le1$,
\[
 |t^{z-1}E(-t)|\le t^{a-1},\qquad e_0(\varepsilon)=\varepsilon^a/a.
\]
Choose an integer $m\ge\max(0,b-1)$.  On $t\ge1$, the absolute
integrand is bounded by $t^mE(-t)$.  Repeated finite integration by parts
on $[T,V]$, followed by the factorial bound for the exponential, gives
\[
 e_\infty(T)=E(-T)m!\sum_{j=0}^m\frac{T^j}{j!}.
\]
This expression tends to zero: for each fixed power, the exponential
series has a higher power as a positive term, giving a rational decay
bound, or repeated doubling of $T$ gives an effective search.  Computed
bounds such as $e_0,e_\infty$ are replaced by rational upper enclosures
when a numerical error budget is chosen.

On each $[\varepsilon,T]$, powers and exponentials have already proved
value, derivative, and continuity estimates, so the middle integral is
an ordinary rectangle computation.  Gamma is therefore constructed for
all these inputs, not introduced by assuming its improper integral exists.
Its standard normalization is recorded in
\href{https://dlmf.nist.gov/5.2}{DLMF, Section 5.2}.

\section{All parameter derivatives, by one envelope}
Fix $z_0$, and choose rational $0<\rho<a\le\Re z_0\le b$.
Expanding $E(wL(t))$ gives
\[
 a_n(t)=t^{z_0-1}E(-t)\frac{L(t)^n}{n!}.
\]
The positive exponential series proves
\[
 \frac{|L(t)|^n}{n!}\le\rho^{-n}
 \begin{cases}t^{-\rho},&0<t\le1,\\t^\rho,&t\ge1.\end{cases}
\]
Thus an integrable coefficient envelope is
\[
 M(t)=\begin{cases}t^{a-\rho-1}E(-t),&t\le1,\\
                   t^{b+\rho-1}E(-t),&t\ge1.
       \end{cases}
\]
Choose an integer $m\ge\max(0,b+\rho-1)$.  The preceding tail computations
apply to $M$, and $\int M\le B=1/(a-\rho)+m!$ is a rational bound.
The parameter-majorant argument supplies a gamma chart of radius $\rho$:
\[
 \Gamma^{(n)}(z_0)=\int_0^\infty t^{z_0-1}E(-t)L(t)^n\,dt,
 \qquad |\Gamma^{(n)}(z_0)|\le n!B\rho^{-n}.
\]
The integral on the right is a computed quantity with its own tails.
Every finite derivative order is now supported; no differentiation of an
uncontrolled infinite integral has occurred.

\section{Recurrence and continuation through known poles}
Integrate the derivative of $t^zE(-t)$ on a truncated interval.  The two
endpoint terms tend to zero by the same estimates, proving
\[
 \Gamma(z+1)=z\Gamma(z),\qquad \Gamma(1)=1.
\]
For a specific $z$ outside the nonpositive integers, choose $N$ so that
$\Re(z+N)>0$ and certify separation of each factor in
\[
 \boxed{\Gamma(z)=\frac{\Gamma(z+N)}{z(z+1)\cdots(z+N-1)}.}
\]
Any two allowable choices agree by the recurrence.  This is an actual
continuation algorithm, since shifting into the right half-plane supplies
new integral bounds, not just a recentering of an old disk.

At $z=-k$ the scalar value is replaced by its pole description.  Isolating
the factor $z+k$ in the denominator and using $\Gamma(1)=1$ gives residue
$(-1)^k/k!$.  The remaining analytic part can be computed by the local
series arithmetic.  For computed inputs, the required separation from
these poles must be supplied; no universal equality decision is asserted.
For positive real arguments gamma is positively separated from zero:
a lower rectangle on $[1,2]$ suffices.  Division by gamma elsewhere
requires its own separation proof; zero-freeness is not assumed here.

\section{Beta and a Gaussian normalization from the circle}
For positive real parameters $a,b$, define
\[
 B(a,b)=\int_0^1 t^{a-1}(1-t)^{b-1}\,dt.
\]
On $(0,1/2]$ the second factor has an explicit rational upper bound;
therefore the left tail is at most a constant times $\varepsilon^a/a$.
Reflection gives the right-tail bound.  The middle integral is bounded
quadrature.  These estimates work uniformly on separated parameter boxes.

\begin{theorem}[Beta--gamma comparison]
\label{thm:beta-gamma}
These independent computations satisfy
\[
 B(a,b)\Gamma(a+b)=\Gamma(a)\Gamma(b).
\]
\end{theorem}
\begin{proof}
Multiply the truncated gamma integrals and interpret the product as finite
rectangular sums in the positive quadrant.  Excluding small coordinate
strips and far tails changes it by the products of the one-dimensional
bounds already proved.  On each remaining bounded region the integrand
is bounded and uniformly continuous.  Triangulate and change the order of
integration; cells meeting the finitely many straight boundaries have
total area tending to zero with the mesh.  Thus the product is equivalently
computed by
\[
 \int_0^\infty E(-r)\int_0^r u^{a-1}(r-u)^{b-1}\,du\,dr.
\]
Here $r=u+v$ is just a shear with determinant $1$, verified on finite
parallelograms; deleting a diagonal tail $r>R$ is controlled by the union
of $u>R/2$ and $v>R/2$.  On each positive $r$ use the one-dimensional
substitution $u=rt$, first away from the two endpoints and then with
the beta endpoint bounds.  The inner integral is $r^{a+b-1}B(a,b)$.
Its remaining endpoint and infinity bounds are the gamma bounds for
$a+b$, proving the identity by finite error comparisons.
\end{proof}

At $a=b=1/2$, use $t=x^2$ and then the rational circle chart
$x=2u/(1+u^2)$, on truncated intervals.  Cancellation gives
\[
 B(1/2,1/2)=4\int_0^1\frac{du}{1+u^2}=\pi.
\]
The endpoint bounds justify both limiting substitutions.  Positivity and
the beta--gamma identity give $\Gamma(1/2)=\sqrt\pi$.  A final
substitution $t=x^2$ therefore proves
\[
 \int_0^\infty E(-x^2)\,dx=\frac{\sqrt\pi}{2}.
\]
The Gaussian normalization has been connected to the earlier circle
computation, without assuming a general two-dimensional change-of-variables
theorem.  Standard beta notation is in
\href{https://dlmf.nist.gov/5.12}{DLMF, Section 5.12}.

\section{The error function: convergent and asymptotic computations}
Define the entire local-series computation
\[
 \operatorname{erf}(z)=\frac2{\sqrt\pi}
       \sum_{n=0}^\infty\frac{(-1)^nz^{2n+1}}{n!(2n+1)}.
\]
On $|z|\le r$, consecutive absolute terms have ratio at most
$r^2/(n+1)$.  After that ratio is at most $1/2$, twice the first omitted
term bounds the tail.  Differentiated tails are obtained in the same way.
Finite integration of the exponential prefixes identifies the function
on the real axis as $2\int_0^x E(-t^2)dt/\sqrt\pi$.

For $x>0$, let $I_k(x)=\int_x^\infty E(-t^2)t^{-2k}dt$.
Finite integration by parts and the vanishing far endpoint give
\[
 I_k(x)=\frac{E(-x^2)}{2x^{2k+1}}-\frac{2k+1}{2}I_{k+1}(x).
\]
Repeated substitution proves, for every finite $N\ge1$,
\[
 I_0(x)=\frac{E(-x^2)}{2x}
   \sum_{k=0}^{N-1}\frac{(-1)^k(2k-1)!!}{(2x^2)^k}+R_N(x),
\]
where $(-1)!!=1$ and
\[
 0\le(-1)^NR_N(x)
 \le\frac{E(-x^2)}{2x}\frac{(2N-1)!!}{(2x^2)^N}.
\]
Indeed $R_N=(-1)^N(2N-1)!!\,I_N/2^N$, and the first integration-by-parts
inequality bounds $I_N$ by $E(-x^2)/(2x^{2N+1})$.

This gives the alternative computation
$1-\operatorname{erf}(x)=2I_0(x)/\sqrt\pi$.  It also illustrates an
important distinction: the asymptotic expression has a proved error for
each \emph{finite} $N$, but its terms eventually grow at fixed $x$.
It is not a convergent infinite series at that input.  Use it when its
finite bound meets the requested precision; otherwise use bounded
quadrature with tails or the convergent series.  The standard asymptotic
formula appears in \href{https://dlmf.nist.gov/7.12}{DLMF, Section 7.12}.
